\documentclass[11pt]{article}
\usepackage{amsmath,amssymb,amsthm}
\usepackage[margin=1in]{geometry}
\newtheorem{lemma}{Lemma}
\newtheorem{theorem}[lemma]{Theorem}
\title{Combining the accepted dispositions for one original marking domain}
\author{Job20 coverage proposal}
\date{}
\begin{document}
\maketitle

These notes give the logical union argument needed for a mixed positive and
negative whole-job proof. They do not acquire an Original domain, replace a
missing proof, or issue an adoption. All model, complete-domain and per-case
premises specified below must be supplied by the exact accepted components.

\paragraph{The domain premise.}
Let $G$ be the named ambient group and let $\mathcal T$ be the entire
Original set of candidate ordered four-mark completions for one fixed prefix.
This set may contain repeated-label tuples, including the indexed candidate
whose added mark equals the first mark. Let $\mathcal A\subseteq\mathcal T$
be the actual admissible rank4 marking domain, which must be contained in
this completed candidate domain.
Suppose an accepted domain certificate gives a disjoint partition
\[
 \mathcal T=\coprod_{p=1}^{N}\mathcal O_p
\]
and a specified representative $t_p$ in each part. Suppose also that every
member of $\mathcal O_p$ is a label-preserving conjugate of $t_p$ by an
element of $G$. Completeness of the candidate domain and containment of the admissible
subset, disjointness, the correct named group and
faithful frame, and the label/member/transporter contract are parts of this
premise. A list of representatives alone does not supply it.

\begin{lemma}\label{lem:transport}
Each of the following properties of an ordered marked tuple is preserved by
simultaneous conjugation in $G$:
\begin{enumerate}
\item its marks generate $G$;
\item it satisfies the intersection property for every pair of label subsets;
\item its marks are distinct nonidentity involutions and satisfy the specified
string commutation relations.
\end{enumerate}
Their negations are preserved as well.
\end{lemma}

\begin{proof}
For a tuple $(g_i)_{i\in I}$ and $h\in G$, conjugation gives, for every
$J\subseteq I$,
\[
 \langle h^{-1}g_jh:j\in J\rangle
   =h^{-1}\langle g_j:j\in J\rangle h.
\]
It maps $G$ to itself and preserves intersections of subgroups. Consequently
the equality between the intersection of the two generated subgroups and
the subgroup generated by the common labels is preserved in both
directions. Conjugation is injective on the labels' values and also preserves identity,
orders and commutators.
These statements prove the claims and their converses. This is the same
label-preserving transport argument already used for the accepted saved-row
properties; it is not a new group computation.
\end{proof}

\begin{theorem}\label{thm:union}
Assume the domain premise. Let $P_N,P_I,P_C$ be disjoint subsets partitioning
$\{1,\ldots,N\}$. Suppose the accepted per-case results prove that:
\begin{enumerate}
\item for every $p\in P_N$, the marks of $t_p$ do not generate $G$;
\item for every $p\in P_I$, the tuple $t_p$ fails the full intersection
property;
\item for every $p\in P_C$, the tuple $t_p$ generates $G$, has the full
intersection property, and has distinct nonidentity involutions and all required string
relations, so is an actual admissible rank4 marking.
\end{enumerate}
Then the entire Original candidate domain is settled relative to these
premises, and in particular every admissible target marking is covered.
Precisely its parts indexed by $P_C$ consist of the required marked string
C-group representations, while no part indexed by $P_N\cup P_I$ does.
\end{theorem}

\begin{proof}
Every Original tuple lies in exactly one part $\mathcal O_p$ by the domain
premise, and has the same relevant properties as $t_p$ by
Lemma~\ref{lem:transport}. If $p\in P_N$, generation fails; if $p\in P_I$,
the required intersection property fails. Either failure excludes a
representation of the required kind. If $p\in P_C$, all required properties
hold. The negative alternatives apply whether or not the candidate was itself
admissible. The positive premise supplies admissibility separately. Since
$\mathcal A\subseteq\mathcal T$, every actual target marking is covered.
These alternatives cover all parts, and their disjointness prevents
counting an Original orbit position twice.
\end{proof}

\paragraph{Recorded job20 partition.}
The metadata reconnaissance joins the final selector, exact edition
78132 bytes with SHA256
\texttt{409f7d5aa45d9c0274374d98c06b04fc241f40e1f08c322db61fa0762cdb9b63},
to its47 accepted source adoptions. Including the original173 conditional
nongeneration positions gives a disjoint positional partition
\[
 N=3779,\qquad |P_N|=232,\quad |P_I|=3542,\quad |P_C|=5,
\]
with
\[
 P_C=\{259,877,1033,1430,2739\}.
\]
The full positional lists and source-field lineage are retained in
\texttt{reconnaissance/partition\_metadata.json}. These are counts of
Original orbit positions, not raw member counts or equivalence classes
under Aut, reversal or duality. The later selector field mentioning only
1430 and2739 is not the complete positive census.

The original173 positions are explicitly identified as conditional
nongeneration in the separately adopted original disposition projection.
They are not inferred from the unrelated all-negative job24 result.
At the earlier small-stabilizer stage, positions877,1033 and1099 had only
generation results and remained unpaid. The positive dispositions877 and1033
and the intersection failure1099 enter through their later accepted results.
Proper-subspace certificates, failed intersection certificates, generation
alone and positive full intersection certificates remain different kinds
of mathematical input.

\paragraph{Original input closure still required.}
Before applying Theorem~\ref{thm:union} to this recorded partition, a whole-job
review must join every family to the same Original prefix20/rank4 domain,
the same14763-member candidate catalogue and3779-row partition, with every
admissible target marking contained in it, and the same exact model,
and its accepted full member/transporter action. The known candidate
certificate pair is the Saved body5a1ae429 and producer body72abfbb4, with
the exact prefix05e8a676, faithful framea8746058 and attributed model bridge
a754edb5. The old context Source06943fa2 checks the complete partition and
the member roles; the present proposal does not execute that Source or
recompute its scientific facts.

The indexed catalogue includes index1=A0, so it can give a repeated mark.
No assertion that all3779 rows already have four distinct admissible marks
is made. Distinctness, nonidentity, involutions, string relations, generation
and the full intersection property must be owned specifically for each of
the five positive cases.

All conditional model, A=G, class, natural-module, faithful-frame,
normalization, maximal-subgroup and transport premises in the accepted
families must remain explicit. A source acceptance saying that those inputs
are borrowed is not a fresh replay of their ancestors. The large inherited
family acceptance26a42fbd is above the4MiB per-file budget and must not be
quietly read through a new bounded metadata entry. Its already accepted
premises may instead be imported through a precise reviewed component, with
the ancestor descriptor and absence of a new rehash stated plainly.

\paragraph{Adoption and remaining scope.}
Only after a distinct complete-domain/family-scope review and Root ownership
of the resulting mixed proof does a new metadata registration become
appropriate. The job24 all-negative proof, final2352 case,173+3606 negative
split and old Root builder API cannot be relabelled as a mixed job20 proof.
The ordered pending-job list would then lose precisely job$[20,4]$; its
global registration must use fresh protected state and retain every other
job and every unrelated field.

Neither the positional union nor the whole-job conclusion settles the
remaining prefixes, the whole field, rank4 output assembly, or final
Aut/reversal/duality equivalences. There is no whole-manuscript, conjecture,
Lean or compression-measurement credit from this coverage proposal.

\end{document}
